
\subsection{Particle Physics}
\subsubsection{Gamma Matrix Identities}
	\begin{minipage}{0.45\textwidth}
		\begin{align*}
			\gamma^{0\dagger}&=\gamma^0
			\\
			\gamma^{j\dagger}&=-\gamma^j\quad\qty(j=1,2,3)
			\\
			\gamma^0\gamma^{\mu\dagger}\gamma^0&=\gamma^\mu
			\\
			\gamma^{0}\gamma^{0}&=\mathbb{I}
			\\
			\gamma^\mu\gamma_\mu&=4
			\\
			\gamma^\mu\gamma_\nu\gamma_\mu&=-2\gamma_\nu
			\\
			\gamma^\mu a_\mu&=\gamma_\mu a^\mu=\slsh{a}
			\\
			\psi^\dagger\gamma^0&=\bar{\psi}
		\end{align*}
	\end{minipage}
	\begin{minipage}{0.48\textwidth}
		\begin{align*}
			\qty{\gamma_\mu,\gamma_\nu}&=\gamma_\mu\gamma_\nu+\gamma_\nu\gamma_\mu=2g_{\mu\nu}
			\\
			\gamma^\mu\gamma_\mu&=4
			\\
			\qty[\gamma^\mu,a_\mu]&=0
			\\
			\gamma^\mu a_\mu&=\gamma_\mu a^\mu=\slsh{a}
			\\
			\Tr(\mathbb{I})&=4
			\\
			\Tr(\gamma_\mu\gamma_\nu)&=\Tr(\gamma^\mu\gamma^\nu)=4
			\\
			\Tr(\gamma^\mu\gamma^\nu\gamma^\rho\gamma^\sigma)&=4(g^{\mu\nu}g^{\rho\sigma}-g^{\mu\rho}g^{\nu\sigma}+g^{\mu\sigma}g^{\nu\rho})
		\end{align*}
	\end{minipage}
	\\[0.25 cm]
	Gamma 5 matrix:
	\\
	\begin{minipage}{0.45\textwidth}
		\begin{align*}
			\gamma_5^\dagger&=\gamma_5
			\\
			\gamma_{5}\gamma_{5}&=\mathbb{I}
			\\
			\qty{\gamma_5,\gamma_\nu}&=0
			\\
			\Tr(\gamma_5)&=0
			\\
			\Tr(\gamma_\mu\gamma_5)&=0
			\\
			\Tr(\gamma_\mu\gamma_\nu\gamma_5)&=0
			\\
			\Tr(\gamma_\mu\gamma_\nu\gamma_\rho\gamma_5)&=0
			\\
			\Tr(\gamma_\mu\gamma_\nu\gamma_\rho\gamma_\sigma\gamma_5)&=4i\epsilon_{\mu\nu\rho\sigma}
		\end{align*}
	\end{minipage}
	\begin{minipage}{0.45\textwidth}
		\begin{align*}
			P_L&=\frac{1-\gamma_5}{2}
			\\
			P_R&=\frac{1+\gamma_5}{2}
			\\
			P_L\gamma_\mu&=\gamma_\mu P_R
			\\
			P_L(\gamma^\mu\gamma^\nu\dotsc\text{odd})P_L&=0
			\\
			P_L(\gamma^\mu\gamma^\nu\dotsc\text{even})P_L&=P_L\qty(\gamma^\mu\gamma^\nu\dotsc)
		\end{align*}
	\end{minipage}
	\\[0.25 cm]
	Quick Complex Conjugates:
	\\
	\begin{align*}
		\qty(\bar{u}\gamma^\mu v)^*&=\bar{v}\gamma^\mu u 
		\\
		\qty(\bar{u}\gamma^\mu P_Lv)^*&=\bar{v}\gamma^\mu P_L u
	\end{align*}
\newpage
\subsubsection{Mandelstam Variables}
	\tikzfeynmanset{every momentum/arrow shorten/.append style={0.7}}
	\begin{align*}
		&
		\feynmandiagram[baseline=(v1.base), horizontal=v1 to v2]{
		a -- [fermion, momentum'=$p_a$] v1 -- [fermion, reversed momentum=$p_b$] b,
		v1 -- [photon, edge label=$p_a+p_b$] v2 -- [fermion, momentum'=$k_1$] c,
		d -- [fermion, momentum=$k_2$] v2,
		};
		&
		&\feynmandiagram[baseline=(v2.base), vertical=v1 to v2]{
			a -- [fermion, momentum'=$p_b$] v1 -- [fermion, momentum=$k_1$] b,
			v1 -- [photon, edge label=$p_a-k_1$] v2 -- [fermion, momentum'=$k_1$] c,
			d -- [fermion, momentum=$p_b$] v2,
			};
		&
		&
		\begin{tikzpicture}
			\begin{feynman}
			\vertex (a) {\(\mu^{-}\)};
			\vertex [right=of a] (b);
			\vertex [above right=of b] (f1) {\(\nu_{\mu}\)};
			\vertex [below right=of b] (c);
			\vertex [above right=of c] (f2) {\(\overline \nu_{e}\)};
			\vertex [below right=of c] (f3) {\(e^{-}\)};
			\diagram* {
			(a) -- [fermion] (b) -- [fermion] (f1),
			(b) -- [boson, edge label'=\(W^{-}\)] (c),
			(c) -- [anti fermion] (f2),
			(c) -- [fermion] (f3),
			};
			\end{feynman}
			\end{tikzpicture}
		\\[0.25cm]
		s&=(p_a+p_b)^2c^2=(k_1+k_2)^2c^2
		&
		t&=(p_a+p_b)^2c^2=(k_1+k_2)^2c^2
	\end{align*}
\newpage
\subsubsection{Feynman Rules}
	\tikzfeynmanset{
		every vertex={red},
		every particle={blue},
		every blob={draw=black, pattern color=black},
		}
	\begin{align*}
		%Incoming anti/fermions
		\feynmandiagram[small, baseline=(v1.base), horizontal=v1 to v2]
		{
			v1 [color=black, particle=$f$] -- [fermion] v2[blob]-- [draw=none] v3 [particle={\phantom{$f$}}, color=white]
		};
		&\quad\longleftrightarrow\quad
		u(p,s)
		\\
		\feynmandiagram[small, baseline=(v1.base), horizontal=v1 to v2]
		{
			v1 [color=black, particle=$f$] -- [anti fermion] v2[blob]-- [draw=none] v3 [particle={\phantom{$f$}}, color=white]
		};
		&\quad\longleftrightarrow\quad
		\overline{v}(p,s)
		\\
		% Outgoing anti/fermions
		\feynmandiagram[small, baseline=(v1.base), horizontal=v1 to v2]
		{
			v3 [draw=none] -- [draw=none] v1 [blob] -- [fermion] v2 [color=black, particle=$f$]
		};
		&\quad\longleftrightarrow\quad
		\overline{u}(p,s)
		\\
		\feynmandiagram[small, baseline=(v1.base), horizontal=v1 to v2]
		{
			v3 [draw=none] -- [draw=none] v1 [blob] -- [anti fermion] v2 [color=black, particle=$f$]
		};
		&\quad\longleftrightarrow\quad
		v(p,s)
		\\
		% Photon fermion interaction
		\feynmandiagram[small, baseline=(v1.base), horizontal=v1 to v2]{
		v1 [particle=$\mu$, color=black] -- [photon] v2,
		[particle=$f$] a -- [fermion] v2 -- [fermion] b [particle=$f$],
		};
		&\quad\longleftrightarrow\quad
		-iQ_fe\gamma^\mu
		\\
	\end{align*}